\subsection{Proof of Proposition \ref{prop:hardness}}
We reduce the $2^n$-tiling problem to the non-emptiness problem for \mbox{B-RASP}.
To this end, we use the following encoding of the function $\tau$ 
as a word over the alphabet $\Sigma := T \cup \{0,1,\#\}$.
We define $\mathit{enc}_\tau \colon \{1,\dots,2^n\} \times \{1,\dots,m\} \to \Sigma^*$ such that
\begin{equation*}\label{eq:enc}
\mathit{enc}_\tau(i,j) := \langle i-1 \rangle \tau(i,j) \#
\end{equation*}
for all $i \in [1,2^n]$ and $j \in [1,m]$, where $\langle i-1 \rangle$ denotes the binary encoding of $i-1$ with $n$ bits
and most significant bit first.
Then 
\[\mathit{enc}(\tau) := \mathit{enc}_\tau(1,1) \dots \mathit{enc}_\tau(2^n,1) \mathit{enc}_\tau(2,1) \dots \mathit{enc}_\tau(m,2^n).\]
We construct a B-RASP program that accepts $\mathit{enc}(\tau)$ if and only if $\tau$ satisfies the conditions above.

Let $n >0$, a finite set $T$ of tiles, and $t_\mathit{fin} \in T$ be given.
The B-RASP program first checks whether the input is a word from
$(\{0,1\}^n T \#)^*$ using the following Boolean vectors:
\begin{align*}
A_T(i) :=&\ \maxatt_j[j<i,1]\ \bigvee_{t \in T} Q_t(j) : 0\\
A_{C,1}(i) :=&\ \maxatt_j[j<i,1]\ Q_0(j) \vee Q_1(j) : 0\\
A_{C,k}(i) :=&\ \maxatt_j[j<i,1]\ A_{C,k-1}(j) : 0 \qquad \text{for } k = 2,\dots,n\\
A_{\#,1}(i) :=&\ \maxatt_j[j<i,1]\ Q_\#(j) : 1\\
A_{\#,k}(i) :=&\ \maxatt_j[j<i,1]\ A_{\#,k-1}(j) : 1 \qquad \text{for } k = 2,\dots,n+1\\
A_{\mathit{enc}}(i) :=&\ \big(Q_\#(i) \to A_T(i)\big) \wedge
\Big(\big(\bigvee_{t \in T} Q_t(i)\big) \to \big(\bigwedge_{k=1}^n A_{C,k}(i)\big) \wedge A_{\#,n+1}(i)\Big)
\end{align*}
We use the vector
\[A(i) :=\ \maxatt_j [j<i,\neg A_\mathit{enc}(j)]\ 0 : A_\mathit{enc}(i)\]
to check that $A_\mathit{enc}(i) = 1$ at every position $i$,
which is the case if and only if $A(\ell) = 1$
where $\ell$ is the length of the input.
Note that we still have to check that the symbol at position $\ell$ is $\#$.
But before that, we ensure that for every two consecutive binary numbers separated by $\#$ the encoded value increases by 1
or is set to 0 if $2^n-1$ is reached.
\begin{align*}
C_1(i) :=&\ \maxatt_j[j<i,Q_0(j) \vee Q_1(j)]\ Q_1(j) : 0\\
C_k(i) :=&\ \maxatt_j[j<i,Q_0(j) \vee Q_1(j)]\ C_{k-1}(j) : 0 \qquad \text{for } k=2,\dots,n\\
C_{+1}(i) :=&\ \maxatt_j[j<i,Q_\#(j)]\ \bigvee_{k=1}^n \big(\bigwedge_{r=1}^{k-1} \neg C_r(i) \wedge C_r(j)\big) \wedge
C_k(i) \wedge \neg C_k(j)\ \wedge\\ &\ \big(\bigwedge_{r=k+1}^{n} C_r(i) \leftrightarrow C_r(j)\big) : 0\\
C_{1 \to 0}(i) :=&\ \maxatt_j[j<i,Q_\#(j)]\ \bigwedge_{k=1}^n \neg C_k(i) \wedge C_k(j) : \bigwedge_{k=1}^n \neg C_k(i)\\
C(i) :=&\ \maxatt_j[j<i,Q_\#(j) \wedge \neg C_{1 \to 0}(j) \wedge \neg C_{+1}(j)]\ 0 : C_{1 \to 0}(i) \wedge C_{+1}(i)
\end{align*} 
Now, $C(\ell) = 1$ if and only if the binary numbers are as required.

Next, we check that the input ends with $1^n t_\mathit{fin} \#$.
\begin{align*}
B_{t}(i) :=&\ \maxatt_j[j<i,\bigvee_{t' \in T} Q_{t'}(j)]\ Q_{t}(j) : 0 \qquad \text{for all } t \in T\\
F(i) :=&\ Q_\#(i) \wedge B_{t_\mathit{fin}}(i) \wedge \bigwedge_{k=1}^n C_k(i)
\end{align*}
Then $F(\ell) = 1$ if and only if the input ends with $1^n t_\mathit{fin} \#$.

We continue by verifying conditions \ref{itm:bot-top} and \ref{itm:first-last} of $\tau$. 
\begin{align*}
E_\bot(i) :=&\ \maxatt_j[j<i,Q_\#(j) \wedge \bigwedge_{k=1}^n C_k(i) \leftrightarrow C_k(j)]\ 1 : 
\bigvee_{t \in T \colon \tdown(t) = 0} B_t(i)\\
E_\top(i) :=&\ \maxatt_j[j<i,Q_\#(j) \wedge 
\big((\bigvee_{t \in T \colon \tup(t) \ne 0} B_t(j)) \vee \bigwedge_{k=1}^n \neg C_k(j)\big)]\ 
\big(\bigvee_{t \in T \colon \tup(t) = 0} B_t(j)\big)\ \wedge\\&\ \big(\bigvee_{t \in T \colon \tup(t) = 0} B_t(i)\big) : 0\\
E_\vdash(i) :=&\ \big(\bigwedge_{k=1}^n \neg C_k(i)\big) \to \big(\bigvee_{t \in T \colon \tleft(t) = 0} B_t(i)\big)\\
E_\dashv(i) :=&\ \big(\bigwedge_{k=1}^n C_k(i)\big) \to \big(\bigvee_{t \in T \colon \tright(t) = 0} B_t(i)\big)\\
E(i) :=&\ \maxatt_j[j<i, Q_\#(j) \wedge \neg(E_\bot(j) \wedge E_\vdash(j) \wedge E_\dashv(j))]\ 0 :
E_\bot(i) \wedge E_\top(i) \wedge E_\vdash(i) \wedge E_\dashv(i)
\end{align*}
Now, conditions \ref{itm:bot-top} and \ref{itm:first-last} hold if and only if $E(\ell) = 1$.

Finally, we ensure that conditions \ref{itm:right-left} and \ref{itm:up-down} are satisfied.
\begin{align*}
M_\downarrow(i) :=&\ \maxatt_j[j<i,Q_\#(j) \wedge \bigwedge_{k=1}^n C_k(i) \leftrightarrow C_k(j)]\ 
\bigvee_{t,t' \in T \colon \tdown(t) = \tup(t')} B_t(i) \wedge B_{t'}(j) : 1\\
M_\leftarrow(i) :=&\ \maxatt_j[j<i,Q_\#(j)]\ \big(\bigvee_{k=1}^n C_k(i)\big) \to 
\big(\bigvee_{t,t' \in T \colon \tleft(t) = \tright(t')} B_t(i) \wedge B_{t'}(j)\big) : 1\\
M(i) :=&\ \maxatt_j[j<i, Q_\#(j) \wedge \neg(M_\downarrow(j) \wedge M_\leftarrow(j))]\ 0 : M_\downarrow(i) \wedge M_\leftarrow(i)
\end{align*}
Then $M(\ell) = 1$ if and only if conditions \ref{itm:right-left} and \ref{itm:up-down} hold.

Thus, if we define the output vector to be the conjunction
\[Y(i) := A(i) \wedge C(i) \wedge F(i) \wedge E(i) \wedge M(i)\]
and say that the B-RASP program accepts if and only if $Y(\ell) = 1$,
then the B-RASP program recognizes the set of all $\mathit{enc}(\tau)$ where $\tau$ satisfies the conditions above.
Hence, the language recognized by the B-RASP program is non-empty if and only if the $2^n$-tiling problem has a solution.


\subsection{Proof of Lemma \ref{lem:equality}}
As affine transformations $A$ and $B$ we use the identity, i.e.,
\[S(\bm v_i,\bm v_j) := \langle \bm v_i, \bm v_j \rangle = \sum_{r=1}^d \big(b_{i,r} b_{j,r} + (1-b_{i,r})(1-b_{j,r})\big)\]
which is equal to $|\{r \in [1,d] \mid b_{i,r} = b_{j,r}\}|$ since
\[b_{i,r} b_{j,r} + (1-b_{i,r})(1-b_{j,r}) = \begin{cases}
1, &\text{if } b_{i,r} = b_{j,r}\\
0, &\text{otherwise.}
\end{cases}\]
Thus, the score is maximized (equal to $d$) if $b_{i,r} = b_{j,r}$ for all $r \in [1,d]$.


\subsection{Proof of Proposition \ref{prop:UHAT-LTL}}
Let $\mathcal{T}$ be a UHAT that recognizes a language $L \subseteq \Sigma^*$ and 
$F$ be a set of binary representations of rational numbers
that may occur during the computation of $\mathcal{T}$ from \cref{lem:UHAT-values}. 
Our goal is to define for the $\ell$-th layer of $\mathcal{T}$ and every vector $\bm v \in F^s$,
where $s$ is the output dimension of layer $\ell$,
an LTL formula $\varphi_{\bm v}^\ell$ such that if $\mathcal{T}$ is applied on input $w \in \Sigma$,
then the $\ell$-th layer outputs at position $i \in [1,|w|]$ the vector $\bm v$ if and only if
$w,i \models \varphi_{\bm v}^\ell$.
We define this formula inductively on the layer number $\ell$.
Let $\emb \colon \Sigma \to (\Q^d)^*$ be the token embedding of $\mathcal{T}$.
For all $\bm v \in F^d$ let
\[\varphi_{\bm v}^0 := \begin{cases}
\bigvee_{a \in \emb^{-1}(\bm v)} Q_a, &\text{if } \emb^{-1}(\bm v) \ne \emptyset\\
\bot, &\text{otherwise.}
\end{cases}\]
We now define the formula for layer $\ell+1$.
In case of a ReLU layer of width $r$, that applies ReLU to the $k$-th coordinate, we can simply define
\[\varphi_{\bm v}^{\ell+1} := \bigvee_{u \in F \colon \max\{0,u\} = \bm v[k]} \varphi_{(\bm v[1,k-1],u,\bm v[k+1,r])}^\ell\]
for all $\bm v \in F^r$.
If layer $\ell+1$ is an attention layer with strict future masking and rightmost tie-breaking defined by the affine transformation
$C \colon \Q^{2r} \to \Q^s$ and score function $S \colon \Q^{2r} \to \Q^r$, we let
\[\varphi_{\bm v}^{\ell+1} := \bigvee_{\substack{\bm u,\bm a \in F^r \colon\\ C(\bm u,\bm a) = \bm v}} \varphi_{\bm u}^\ell \wedge
\big((\bigvee_{\substack{\bm b \in F^r \colon\\ S(\bm u,\bm b) < S(\bm u,\bm a)}} \varphi_{\bm b}^\ell) \since
(\varphi_{\bm a}^\ell \wedge \neg\past \bigvee_{\substack{\bm b \in F^r \colon\\ S(\bm u,\bm b) > S(\bm u,\bm a)}} 
\varphi_{\bm b}^\ell)\big)\]
for all $\bm v \in F^s$.
To account for the special case, where the set of unmasked positions is empty, we take the disjunction of the previous formula and
$(\neg\past\top) \wedge \bigvee_{\bm u \in F^r \colon C(\bm u,\bm 0) = \bm v} \varphi_{\bm u}^\ell$.
We omit this special case in the following.
If the layer uses leftmost tie-breaking, we adapt the formula as follows:
\[\varphi_{\bm v}^{\ell+1} := \bigvee_{\substack{\bm u,\bm a \in F^r \colon\\ C(\bm u,\bm a) = \bm v}} \varphi_{\bm u}^\ell \wedge
\big(\past (\varphi_{\bm a}^\ell \wedge \neg\past \bigvee_{\substack{\bm b \in F^r \colon\\ S(\bm u,\bm b) \ge S(\bm u,\bm a)}} 
\varphi_{\bm b}^\ell)\big) \wedge
\big(\neg\past \bigvee_{\substack{\bm b \in F^r \colon\\ S(\bm u,\bm b) > S(\bm u,\bm a)}} 
\varphi_{\bm b}^\ell)\big)\]
The case of strict past masking is similar, where we use $\until$ instead of $\since$ and $\future$ instead of $\past$.
If the layer uses no masking and rightmost tie-breaking, we distinguish three situations:
the attention vector is at the current position, the attention vector is strictly to the left of the current position,
or the attention vector is strictly to the right of the current position.
For the situation, where the attention vector is at the current position, we use
\begin{equation}\label{eq:no-masking1}
\bigvee_{\substack{\bm u \in F^r \colon\\ C(\bm u,\bm u) = \bm v}} \varphi_{\bm u}^\ell \wedge
\big(\neg\past\bigvee_{\substack{\bm b \in F^r \colon\\ S(\bm u,\bm b) > S(\bm u,\bm u)}} 
\varphi_{\bm b}^\ell\big) \wedge
\big(\neg\future\bigvee_{\substack{\bm b \in F^r \colon\\ S(\bm u,\bm b) \ge S(\bm u,\bm u)}} 
\varphi_{\bm b}^\ell\big)
\end{equation}
For the situation, where the attention vector is strictly to the left of the current position, we use
\begin{equation}\label{eq:no-masking2}
\begin{split}
\bigvee_{\substack{\bm u,\bm a \in F^r \colon\\ 
C(\bm u,\bm a) = \bm v \wedge S(\bm u,\bm a) > S(\bm u,\bm u)}} &\varphi_{\bm u}^\ell \wedge
\big(\neg\future\bigvee_{\substack{\bm b \in F^r \colon\\ S(\bm u,\bm b) \ge S(\bm u,\bm a)}} 
\varphi_{\bm b}^\ell\big) \\&\wedge
\big((\bigvee_{\substack{\bm b \in F^r \colon\\ S(\bm u,\bm b) < S(\bm u,\bm a)}} \varphi_{\bm b}^\ell) \since
(\varphi_{\bm a}^\ell \wedge \neg\past \bigvee_{\substack{\bm b \in F^r \colon\\ S(\bm u,\bm b) > S(\bm u,\bm a)}} 
\varphi_{\bm b}^\ell)\big).
\end{split}
\end{equation}
Similarly, for the situation, where the attention vector is strictly to the right of the current position, we use
\begin{equation}\label{eq:no-masking3}
\begin{split}
\bigvee_{\substack{\bm u,\bm a \in F^r \colon\\ 
C(\bm u,\bm a) = \bm v \wedge S(\bm u,\bm a) \ge S(\bm u,\bm u)}} &\varphi_{\bm u}^\ell \wedge
\big(\neg\past\bigvee_{\substack{\bm b \in F^r \colon\\ S(\bm u,\bm b) > S(\bm u,\bm a)}} 
\varphi_{\bm b}^\ell\big) \\&\wedge
\big(\future (\varphi_{\bm a}^\ell \wedge \neg\future \bigvee_{\substack{\bm b \in F^r \colon\\ S(\bm u,\bm b) \ge S(\bm u,\bm a)}} 
\varphi_{\bm b}^\ell)\big) \wedge
\big(\neg\future \bigvee_{\substack{\bm b \in F^r \colon\\ S(\bm u,\bm b) > S(\bm u,\bm a)}} 
\varphi_{\bm b}^\ell)\big).
\end{split}
\end{equation}
Thus, in the case of no masking and rightmost tie-breaking, we define $\varphi_{\bm v}^{\ell+1}$ 
as the disjunction of \cref{eq:no-masking1,eq:no-masking2,eq:no-masking3}.
The case where the layer uses no masking and leftmost tie-breaking is analogous.

Finally, if there are $m$ layers, where the last layer outputs vectors of dimension $s$, 
and $\bm t \in \Q^s$ is the acceptance vector of $\mathcal{T}$, we define the formula
\[\varphi := \bigvee_{\bm v \in F^s \colon \\ \langle \bm t,\bm v\rangle > 0} \varphi_{\bm v}^m.\]
Then $w,k \models \varphi$ if and only if $w \in L$, where $k$ is the output position of $\mathcal{T}$.

It remains to argue that $\varphi$ can be computed in exponential time.
By \cref{lem:UHAT-values}, $|F|$ is exponential in the size of $\mathcal{T}$ 
and every representation in $F$ is of polynomial size.
Moreover, $F$ can be computed in exponential time.
The formulas $\varphi_{\bm v}^{\ell+1}$ at every layer $\ell+1$ of width $r$
depends on $|F|^{O(r)}$ many formulas from layer $\ell$.
Moreover, $\varphi_{\bm v}^{\ell+1}$ can be computed in time polynomial in $|F|^r \cdot |\mathcal{T}|$,
since we only have to compute affine transformations on vectors from $F^r$,
where each component is of size polynomial in $|\mathcal{T}|$.
The formulas $\varphi_{\bm v}^{m}$ at the last layer $m$ depend on $|F|^{O(r'm)}$ many formulas from layer 0,
where $r'$ is the maximum width of all layers.
Thus, $\varphi_{\bm v}^{m}$ has size exponential in $|\mathcal{T}|$ and can be computed in exponential time.
Therefore, also $\varphi$ can be computed in exponential time.